
\section{Background}

\subsection{Multi-Head Self-Attention}

For input $X \in \mathbb{R}^{B \times L \times d}$:

\[
\begin{aligned}
Q &= X W\_Q, \quad K = X W\_K, \quad V = X W\_V \\
\text{head}\_h &= \mathrm{softmax}\!\left(\frac{Q\_h K\_h^\top}{\sqrt{d\_h}}\right) V\_h \\
\text{MHSA}(X) &= \mathrm{Concat}(\text{head}\_1, \dots, \text{head}\_H) W\_O
\end{aligned}
\tag{3.1}
\]

\subsection{Rotary Position Embedding (RoPE)}

RoPE encodes position by rotating each 2D subspace of $Q, K$:

\[
\tilde{q}\_i = q\_i \cos(\theta\_i) + \mathrm{rot}(q\_i) \sin(\theta\_i)
\tag{3.2}
\]

where the rotation operator acts on the first half vs second half:

\[
\mathrm{rot}(x) = \begin{pmatrix} -x\_{d\_h/2:} \\ x\_{:d\_h/2} \end{pmatrix}
\tag{3.3}
\]

and $\theta_i$ is a position-dependent angle:

\[
\theta\_i^{(k)} = i \cdot \omega\_k, \qquad
\omega\_k = \mathrm{base}^{-2k/d\_h}, \quad k = 0, \dots, d\_h/2 - 1
\tag{3.4}
\]

\textbf{Effect on attention.} The score between $q_i$ and $k_j$ depends only on 
$i - j$:

\[
\tilde{q}\_i^\top \tilde{k}\_j = f(q\_i, k\_j, i - j)
\tag{3.5}
\]

\subsection{Partial-Dimension RoPE}

Not all models apply RoPE to the full head dimension. For models with 
partial rotation (e.g., Qwen with \texttt{rotary\_dim \textless{} head\_dim}):

\[
\tilde{q}\_i = \begin{pmatrix} q\_i^{[:r]} \cos(\theta\_i) + \mathrm{rot}(q\_i^{[:r]}) \sin(\theta\_i) \\ q\_i^{[r:]} \end{pmatrix}
\tag{3.6}
\]

where $r$ is the rotary dimension. UnifiedRoPE supports this via 
\texttt{rotary\_dim = 2 * (head\_dim // 2)}.

\subsection{Grouped-Query Attention (GQA)}

$H_{kv} < H$ KV heads, repeated $H / H_{kv}$ times before attention:

\[
K\_\text{expanded} = \mathrm{repeat\_kv}(K, H / H\_{kv})
\tag{3.7}
\]

UnifiedRoPE applies \texttt{repeat\_kv} after RoPE.

\subsection{Nested Encoder Structure}

In BERT-style models:

\begin{lstlisting}
BertAttention
+-- self       (BertSelfAttention)
|   +-- query
|   +-- key
|   +-- value
+-- output     (BertSelfOutput)
    +-- dense  # $\leftarrow$ Output projection OUTSIDE
\end{lstlisting}

The output projection \texttt{output.dense} is \textbf{outside} the inner attention. 
When replacing the inner attention, the outer \texttt{output.dense} must still 
apply. UnifiedRoPE returns raw context and lets the outer layer project.

\subsection{Conv1D in GPT-2}

GPT-2 uses HuggingFace's \texttt{Conv1D}, which stores weights as \texttt{(in, out)}:

\begin{lstlisting}
class Conv1D(nn.Module):
    def \_\_init\_\_(self, nf, nx):
        self.weight = nn.Parameter(torch.empty(nx, nf))  # (in, out)
\end{lstlisting}

This is \textbf{transposed} relative to \texttt{nn.Linear}. UnifiedRoPE detects Conv1D 
and transposes during weight transfer.

\subsection{LoRA Adapters}

LoRA (Hu et al., 2021) adds trainable residuals:

\[
\Delta W = \frac{\alpha}{r} B A
\tag{3.8}
\]

with $B = 0$ at initialization. UnifiedRoPE's \texttt{LoRALinear} wraps this.

\subsection{Precomputed Caches}

RoPE angles $\theta_i$ are precomputed once:

\begin{lstlisting}
inv\_freq = 1.0 / (base ** (arange(0, rotary\_dim, 2) / rotary\_dim))
t = arange(seq\_len)
freqs = einsum("i,j-\textgreater{}ij", t, inv\_freq)
emb = cat([freqs, freqs], dim=-1)
cos\_cached = emb.cos()
sin\_cached = emb.sin()
\end{lstlisting}

UnifiedRoPE stores these as \textbf{non-persistent buffers} (\texttt{persistent=False}), 
so they are not saved with the model state dict.

\subsection{Notation Summary}

\begin{lstlisting}

---
\end{lstlisting}